% =====================================================================
%  3 REFERENCIAL TEÓRICO — "Quais conceitos subsidiam o trabalho?"
%  Cada subseção deve terminar ligando-se à seguinte (Dica 2), e todas
%  devem ter relação direta com os resultados.
%  Tabelas/figuras sempre citadas pelo número: "conforme a Tabela 1".
% =====================================================================
\section{Referencial Teórico}\label{sec:referencial}

\escrever{parágrafo de abertura explicando a sequência dos conceitos.}

\subsection{Sensoriamento tátil}
\escrever{tecnologias de transdução (resistiva, capacitiva, piezoelétrica, óptica e magnética) e por que a magnética é adequada à ponta do dedo \cite{dahiya2010tactile}.}

\subsection{Efeito Hall e magnetômetros tridimensionais}
\escrever{princípio do efeito Hall; campo de um ímã permanente e decaimento com a distância; características do TLV493D \cite{infineon2016tlv493d}: faixa, resolução, endereçamento \ac{I2C}.}

\begin{equation}
  B_z(z) = \frac{B_r}{2}\left[\frac{z+L}{\sqrt{R^2+(z+L)^2}} - \frac{z}{\sqrt{R^2+z^2}}\right]
  \label{eq:campo-disco}
\end{equation}
\escrever{explicar cada símbolo da Equação~\eqref{eq:campo-disco} (campo axial de um ímã em disco).}

\subsection{Microestruturas e manufatura aditiva em TPU}
\escrever{microestruturas cut-cell \cite{tozoni2024cutcell}, rigidez controlada pela geometria \cite{panetta2015elastic}, comportamento do \ac{TPU} impresso sob compressão \cite{bates2016honeycombs}.}

\subsection{Calibração por aprendizado de máquina}
\escrever{regressão linear e polinomial; \ac{MLP}; métricas (\ac{RMSE}, R$^2$); validação temporal e vazamento de dados \cite{bergmeir2012crossvalidation}; inferência embarcada \cite{david2021tflm}.}

\subsection{Feedback vibrotátil}
\escrever{mecanorreceptores e percepção de vibração; atuadores \ac{ERM} e \ac{LRA} \cite{culbertson2018haptics}; acionamento por \ac{PWM}.}
